\section{Conclusion}
\label{sec:conclusion}

This study presents an end-to-end approach for active 3D scene reconstruction, reducing the need for manual intervention. 
Specifically, the learning-based policy explores how to reconstruct diverse objects in the training stage and thus can generalize to reconstruct unseen objects in a fully autonomous manner.
Our controller maneuvers in free space and selects the next best view based on a hybrid scene representation which conveys scene coverage status and thus reconstruction progress. 
We show the effectiveness of our approach by generalizing it on multiple datasets. %  including Houses3K, OmniObject3D, Objaverse and Replica
The quantitative and qualitative generalization results on holdout Houses3K test set and cross-domain OmniObject3D, including house category and non-house categories, show that our method outperforms other baselines in terms of reconstruction completeness, efficiency and accuracy. 
Furthermore, the experiment conducted on Objaverse and Replica shows that the policy trained in house-only settings can even generalize to complicated outdoor and indoor scenes.

% \noindent \textbf{Acknowledgements.} This work is supported by Shanghai Artificial Intelligence Laboratory and CUHK Direct Grants (RCFUS) 4055189. We're very grateful to Tao Huang, Zeqi Xiao, Junfeng Long, Mulin Yu, and all reviewers for their insightful advice.
\vspace{4pt}
\section*{Acknowledgements}
This work is supported by Shanghai Artificial Intelligence Laboratory and CUHK Direct Grants (RCFUS) 4055189. We're very grateful to Tao Huang, Zeqi Xiao, Junfeng Long, Mulin Yu, and all reviewers for their insightful advice.

\newpage